\en{In this chapter, we prove some properties of the Weierstraß $\sigma$-function and the Fourier representations of Thm.~\ref{\en{EN}fouriertheorem}.
The proof follows \cite[ch.~18, §1-3]{\en{EN}Lang1987}.}%
\de{In diesem Kapitel beweisen wir einige Eigenschaften der Weierstraß'schen $\sigma$-Funktion und die Fourierdarstellungen in Thm.~\ref{\en{EN}fouriertheorem}. Der Beweis folgt \cite[Kap.~18, §1-3]{\en{EN}Lang1987}.}

\begin{thm}\label{\en{EN}sigmaodd}
\en{The Weierstraß $\sigma$-function is an odd function: $\sigma(-z;L)=-\sigma(z;L)$.}%
\de{Die Weierstraß'sche $\sigma$-Funktion ist ungerade: $\sigma(-z;L)=-\sigma(z;L)$.}
\end{thm}
\begin{proof}
\en{We recall Definition~\ref{\en{EN}defisigma} from page~\pageref{\en{EN}defisigma}:}%
\de{Wir rufen uns Definition~\ref{\en{EN}defisigma} von Seite~\pageref{\en{EN}defisigma} in Erinnerung:}
$$\sigma(z;L):=z\cdot\prod_{\substack{\omega\in L\\\omega\neq 0}}\left\{\left(1-\frac z \omega \right)\cdot\exp\lk\frac z \omega + \frac 1 2 \left(\frac z \omega\right)^2\rk\right\}$$
\en{Here we realize that if $\omega$ runs through the whole lattice $L$, then $-\omega$ does the same. This yields $\sigma(-z;L)=-\sigma(z;L)$, where the additional minus sign comes from the factor $z$ in front of the product sign.}%
\de{Hier erkennen wir, dass mit $\omega$ auch $-\omega$ alle Punkte des Gitters $L$ durchläuft, und dass folglich $\sigma(-z;L)=-\sigma(z;L)$ gilt, wobei das zusätzliche Minuszeichen vom ersten Faktor $z$ stammt.}
\end{proof}

\begin{thm}\label{\en{EN}trafosigma}
\en{When translating $z$ by one of the basic periods of the lattice $L=\mathbb Z \omega_1 + \mathbb Z \omega_2$, the Weierstraß $\sigma$-function transforms as follows: }%
\de{Es gilt das folgende Transformationsverhalten der $\sigma$-Funktion bei Translation um eine der beiden Basisperioden des Gitters $L=\mathbb Z \omega_1 + \mathbb Z \omega_2$:}
$$\sigma(z+\omega_k)=-\exp\lk\eta_k\cdot\left(z+\frac{\omega_k}{2}\right)\rk\cdot \sigma(z)$$
\end{thm}
\begin{proof}
\en{From the Definitions~\ref{\en{EN}defizeta} and~\ref{\en{EN}defetak} of $\zeta$ and the $\eta_k$ we get:}%
\de{Zunächst folgt aus den Definitionen~\ref{\en{EN}defizeta} und~\ref{\en{EN}defetak} für $\zeta$ und die $\eta_k$:}
\begin{align*}
    \frac{d}{dz}\log\lk\frac{\sigma(z+\omega_k)}{\sigma(z)}\rk
    &= \frac{d}{dz}\log\lk\sigma(z+\omega_k)\rk - \frac{d}{dz}\log\lk\sigma(z)\rk\\
    &= \frac{\sigma'(z+\omega_k;L)}{\sigma(z+\omega_k;L)} - \frac{\sigma'(z;L)}{\sigma(z;L)} =\zeta(z+\omega_k;L)-\zeta(z;L) =\eta_k\\
    \Longrightarrow\quad \log\lk\frac{\sigma(z+\omega_k)}{\sigma(z)}\rk &= \eta_k\cdot z + c(\omega_k)\\
    \Longrightarrow\quad \frac{\sigma(z+\omega_k)}{\sigma(z)} &= \exp\lk\eta_k\cdot z + c(\omega_k)\rk = \exp\lk\eta_k\cdot z\rk\cdot\exp\lk c(\omega_k)\rk
\end{align*}
\en{We get the value of $\exp\lk c(\omega_k)\rk$ by setting $z=-\frac{\omega_k}{2}$ and using the fact that $-\frac{\omega_k}{2}\notin L_\tau$ and that $\sigma$ is odd (see Prop.~\ref{\en{EN}sigmaodd}):}%
\de{Den Wert $\exp\lk c(\omega_k)\rk$ bestimmen wir, indem wir $z=-\frac{\omega_k}{2}$ einsetzen. Dabei verwenden wir, dass dieser Wert nicht im Gitter $L_\tau$ liegt und dass $\sigma$ ungerade ist:}\belowdisplayskip=-6pt
\begin{align*}
     \exp\lk-\eta_k\cdot\frac{\omega_k}{2} + c(\omega_k)\rk
     &= \frac{\sigma\lk-\frac{\omega_k}{2}+\omega_k\rk}{\sigma\lk-\frac{\omega_k}{2}\rk}
     = \frac{\sigma\lk\frac{\omega_k}{2}\rk}{\sigma\lk-\frac{\omega_k}{2}\rk}
     = -1\quad\left|~\cdot~\exp\lk\eta_k\cdot\frac{\omega_k}{2}\rk\right.\\
     \Longrightarrow\quad \exp\lk c(\omega_k)\rk &= -\exp\lk\eta_k\cdot\frac{\omega_k}{2}\rk\\
     \Longrightarrow\quad \frac{\sigma(z+\omega_k)}{\sigma(z)}
     &= -\exp\lk\eta_k\cdot z\rk\cdot\exp\lk\eta_k\cdot\frac{\omega_k}{2}\rk
     = -\exp\lk\eta_k\cdot\left(z+\frac{\omega_k}{2}\right)\rk
\end{align*}
\end{proof}



\begin{thm}\label{\en{EN}satzphi}
 \en{We define the function}\de{Wir definieren die Funktion}
 $$\varphi(z;L_\tau):=\exp\lk-\frac {\eta_1}{2}\cdot z^2+i\pi z\rk\cdot\sigma(z;L_\tau)$$
 \en{using the first basic quasiperiod $\eta_1=\eta_1(L_\tau)$. For this function, it holds:}%
 \de{mit Hilfe der ersten Basisquasiperiode $\eta_1=\eta_1(L_\tau)$. Für diese Funktion gilt dann:}
 $$\varphi(z+1;L_\tau)=\varphi(z;L_\tau)\qquad\text{\en{and}\de{und}}\qquad\varphi(z+\tau;L_\tau) = -\exp(2\pi i z)\cdot \varphi(z;L_\tau)$$
\end{thm}
\begin{proof}
\en{We use the transformation formula of the $\sigma$-function from Prop.~\ref{\en{EN}trafosigma}:}%
\de{Wir verwenden das Transformationsverhalten der $\sigma$-Funktion aus Satz~\ref{\en{EN}trafosigma}:}
\begin{align*}
    \varphi(z+1;L_\tau) &= \exp\lk-\frac 1 2 \eta_1\cdot (z+1)^2+i\pi (z+1)\rk\cdot\sigma(z+1;L_\tau)\\
    &= -\exp\lk-\frac {\eta_1}{2}\cdot (z^2+2z+1) +i\pi (z+1)+\eta_1\cdot\left(z+\frac{\omega_1}{2}\right)\rk\cdot\sigma(z;L_\tau)\\
    &= -\exp\lk-\frac {\eta_1}{2}\cdot (2z+1) +i\pi +\eta_1\cdot\left(z+\frac{1}{2}\right)\rk\cdot\varphi(z;L_\tau)\\
    &= -\exp\lk i\pi\rk\cdot\varphi(z;L_\tau) = \varphi(z;L_\tau)
\end{align*}
\en{For the second basic period of $L_\tau$, Prop.~\ref{\en{EN}trafosigma} yields:}%
\de{Und für die andere Basisperiode von $L_\tau$ folgt ebenfalls mit Satz~\ref{\en{EN}trafosigma}:}
\begin{align*}
    \varphi(z+\tau;L_\tau) &= \exp\lk-\frac {\eta_1}{2}\cdot (z+\tau)^2+i\pi (z+\tau)\rk\cdot\sigma(z+\tau;L_\tau)\\
    &= -\exp\lk-\frac {\eta_1}{2}\cdot (z^2+2z\tau+\tau^2) +i\pi (z+\tau)+\eta_2\cdot\left(z+\frac{\omega_2}{2}\right)\rk\cdot\sigma(z;L_\tau)\\
    &= -\exp\lk-\frac {\eta_1}{2}\cdot (2z\tau+\tau^2) +i\pi \tau+\eta_2\cdot\left(z+\frac{\tau}{2}\right)\rk\cdot\varphi(z;L_\tau)\\
    &= -\exp\lk-\eta_1\cdot\tau \left(z+\frac \tau 2\right) +i\pi \tau+\eta_2\cdot\left(z+\frac{\tau}{2}\right)\rk\cdot\varphi(z;L_\tau)
\end{align*}
\en{Then we use Legendre's relation of the lattice $L_\tau$ from Prop.~\ref{\en{EN}legendre}, which reads $\eta_1\cdot\tau = 2\pi i + \eta_2$ (because $\omega_1=1$ and $\omega_2=\tau$). This yields:}%
\de{Jetzt verwenden wir noch die Legendre-Relation des Gitters $L_\tau$ aus Satz~\ref{\en{EN}legendre} in der Form $\eta_1\cdot\tau = 2\pi i + \eta_2$ (setze $\omega_1=1$ und $\omega_2=\tau$ ein) und erhalten:}\belowdisplayskip=0pt
\begin{align*}
    \varphi(z+\tau;L_\tau)
        &= -\exp\lk-(2\pi i + \eta_2)\cdot \left(z+\frac \tau 2\right) +i\pi \tau+\eta_2\cdot\left(z+\frac{\tau}{2}\right)\rk\cdot\varphi(z;L_\tau)\\
        &= -\exp\lk-2\pi i \left(z+\frac \tau 2\right) +i\pi \tau\rk\cdot\varphi(z;L_\tau)= -\exp\lk-2\pi i z\rk\cdot\varphi(z;L_\tau)
\end{align*}
\end{proof}

\begin{thm}\label{\en{EN}fouriersigma}
\en{The Weierstraß $\sigma$-function of the lattice $L_\tau$ admits the following Fourier series expansion with $q_z=e^{2\pi i z}$ and $q_\tau=e^{2\pi i \tau}$ and $\eta_1=\eta_1(L_\tau)$:}%
\de{Die Weierstraß'sche $\sigma$-Funktion zum Gitter $L_\tau$ hat die folgende Fourierentwicklung mit $q_z=e^{2\pi i z}$ und $q_\tau=e^{2\pi i \tau}$ und $\eta_1=\eta_1(L_\tau)$:}
$$\sigma(z;\tau)=\frac 1 {2\pi i}e^{\eta_1 \cdot z^2/2}\cdot(q_z^{1/2}-q_z^{-1/2})\cdot\prod_{n=1}^\infty\frac{(1-q_\tau^n q_z)(1-q_\tau^n / q_z)}{(1-q_\tau^n)^2}$$
\end{thm}
\begin{proof}
\en{First we will prove that for the $\varphi$-function from Prop.~\ref{\en{EN}satzphi} it holds:}%
\de{Wir beweisen zunächst, dass für die in Satz~\ref{\en{EN}satzphi} definierte $\varphi$-Funktion gilt:}
\begin{align}
    \varphi(z;L_\tau)=\frac {q_z-1} {2\pi i}\cdot\prod_{n=1}^\infty\frac{(1-q_\tau^n q_z)(1-q_\tau^n / q_z)}{(1-q_\tau^n)^2} \label{\en{EN}glgphi}
\end{align}
\en{We denote the right hand side of~(\ref{\en{EN}glgphi}) by $g(z;L_\tau)$ and prove that $g(z;L_\tau)=\varphi(z;L_\tau)$:
From $q_{z+1}=e^{2\pi i (z+1)} = e^{2\pi i z}= q_z$ we get $g(z+1;L_\tau)=g(z;L_\tau)$, like with $\varphi(z;L_\tau)$ (cf.~Prop.~\ref{\en{EN}satzphi}).
Then it holds $q_{z+\tau}=q_z\cdot q_\tau$ and thus}%
\de{Die rechte Seite von~(\ref{\en{EN}glgphi}) nennen wir $g(z;L_\tau)$ und beweisen dann $g(z;L_\tau)=\varphi(z;L_\tau)$:
Es gilt $q_{z+1}=e^{2\pi i (z+1)} = e^{2\pi i z}= q_z$ und somit $g(z+1;L_\tau)=g(z;L_\tau)$, genau wie bei $\varphi(z;L_\tau)$ (vgl.~Satz~\ref{\en{EN}satzphi}). Außerdem ist $q_{z+\tau}=q_z\cdot q_\tau$ und somit}
\begin{align*}
    g(z+\tau;L_\tau) &=\frac {q_zq_\tau-1} {2\pi i}\cdot\prod_{n=1}^\infty\frac{(1-q_\tau^{n+1} q_z)(1-q_\tau^{n-1} / q_z)}{(1-q_\tau^n)^2}\\
    &= \frac {q_zq_\tau-1} {2\pi i}\cdot\frac{\left\{\prod_{n=2}^\infty(1-q_\tau^{n} q_z)\right\}\cdot\left\{\prod_{n=0}^\infty(1-q_\tau^n / q_z)\right\}}{\prod_{n=1}^\infty(1-q_\tau^n)^2}\\
    &= \frac {q_zq_\tau-1} {2\pi i}\cdot \frac{1-q_\tau^0/q_z}{1-q_\tau^1q_z}
    \cdot \prod_{n=1}^\infty\frac{(1-q_\tau^{n+1} q_z)(1-q_\tau^{n-1} / q_z)}{(1-q_\tau^n)^2}\\
    &= \frac {q_zq_\tau-1} {2\pi i}\cdot \frac{1-1/q_z}{1-q_\tau q_z}
    \cdot \frac {q_z-1}{q_z-1}\cdot\prod_{n=1}^\infty\frac{(1-q_\tau^{n+1} q_z)(1-q_\tau^{n-1} / q_z)}{(1-q_\tau^n)^2}\\
    &=\frac {q_zq_\tau-1} {q_z-1}\cdot \frac{1-1/q_z}{1-q_\tau q_z} \cdot g(z;L_\tau) =\frac {q_zq_\tau-1}{1-q_\tau q_z}\cdot \frac{q_z^{-1}(q_z-1)} {q_z-1} \cdot g(z;L_\tau)\\
    &= - q_z^{-1}\cdot  g(z;L_\tau) = -\exp(-2\pi i z) \cdot g(z;L_\tau)
\end{align*}
\en{Thus we have shown that $g$ and $\varphi$ act in the same way when transforming $z\rightarrow z+1$ or $z\rightarrow z+\tau$. This shows that $\frac{\varphi}g$ is doubly periodic with period lattice $L_\tau$.}%
\de{Somit haben wir bewiesen, dass sich $g$ und $\varphi$ bei $z\rightarrow z+1$ und bei $z\rightarrow z+\tau$ genau gleich verhalten, dass also $\frac{\varphi}g$ das Gitter $L_\tau$ als Periodengitter hat.}

\en{Now we analyze the zeros of $g(z;L_\tau)$. From the rule of zero product we deduce that $g(z;L_\tau)=0$ iff $q_z=q_\tau^m$ for any $m\in \mathbb Z$.
This yields the condition $e^{2\pi i z} = e^{2\pi i m \tau}$ for the zeros of $g$.
But since the natural exponential function has the complex period $2\pi i$, every $(l,m)\in\mathbb Z^2$ produces a zero of $g$:
$2\pi i \cdot z =  2\pi i \cdot l + 2\pi i \cdot m \tau $ or $z=l + m\tau$. Thus $g(z;L_\tau)$ has zeros of order one for $z\in L_\tau$, like $\sigma(z;L_\tau)$ (cf.~Remark~\ref{\en{EN}bemsigma}) and like $\varphi(z;L_\tau)$ (cf.~Prop.~\ref{\en{EN}satzphi}).
}%
\de{Nun kommen wir zu den Nullstellen von $g(z;L_\tau)$. Aus dem Satz vom Nullprodukt folgt, dass $g(z;L_\tau)=0$ genau dann gilt, wenn $q_z=q_\tau^m$ für ein $m\in \mathbb Z$.
Dies liefert die Gleichung $e^{2\pi i z} = e^{2\pi i m \tau}$. Aufgrund der komplexen Periode $2\pi i$ der $e$-Funktion erhalten wir für alle $(l,m)\in\mathbb Z^2$ eine Lösung:
$2\pi i \cdot z =  2\pi i \cdot l + 2\pi i \cdot m \tau $ bzw.~$z=l + m\tau$. Folglich hat $g(z;L_\tau)$ in allen Punkten des Gitters $L_\tau$ eine einfache Nullstelle, genau wie $\sigma(z;L_\tau)$ (siehe Bemerkung~\ref{\en{EN}bemsigma}) und somit genau wie $\varphi(z;L_\tau)$ (siehe Satz~\ref{\en{EN}satzphi}).}

\en{From its definition in Prop.~\ref{\en{EN}satzphi} we deduce that $\varphi$ has no poles, and thus $\frac{\varphi}{g}$ is an elliptic function without poles and must be constant (first Liouvielle theorem, Prop.~\ref{\en{EN}liouville1}).}%
\de{Weil zudem $\varphi$ aufgrund der Definition in Satz~\ref{\en{EN}satzphi} keine Polstellen hat, ist die elliptische Funktion $\frac{\varphi}{g}$ wegen des ersten Liouville'schen Satzes~\ref{\en{EN}liouville1} konstant.}

\en{Next we calculate this constant value for $z\rightarrow 0$.
There we have $q_z= 1+2\pi i z + O(z^2)$ and thus $g(z;L_\tau)\approx \frac{1+2\pi i z-1}{2\pi i}\cdot \prod_{n=1}^\infty\frac{(1-q_\tau^n)(1-q_\tau^n )}{(1-q_\tau^n)^2}=\frac {2\pi i z}{2\pi i}\cdot 1 = z$.
We deduce the same approximation $\sigma(z;L_\tau)\approx z$ around $z=0$ from the Definition~\ref{\en{EN}defisigma} of $\sigma(z;L_\tau)$.
Thus it holds $\varphi(z;L_\tau)\approx z$ around $z=0$ and we get:}%
\de{Den Wert der Konstante bestimmen wir für $z\rightarrow 0$.
Dort ist $q_z= 1+2\pi i z + O(z^2)$ und somit $g(z;L_\tau)\approx \frac{1+2\pi i z-1}{2\pi i}\cdot \prod_{n=1}^\infty\frac{(1-q_\tau^n)(1-q_\tau^n )}{(1-q_\tau^n)^2}=\frac {2\pi i z}{2\pi i}\cdot 1 = z$.
Die gleiche Näherung $\sigma(z;L_\tau)\approx z$ in der Nähe von $z=0$ erkennt man an der Def.~\ref{\en{EN}defisigma} von $\sigma(z;L_\tau)$.
Somit gilt auch $\varphi(z;L_\tau)\approx z$ in der Nähe von $z=0$ und somit:}
$$\frac{\varphi(z;L_\tau)}{g(z;L_\tau)}=\lim_{z\rightarrow 0}\frac{\varphi(z;L_\tau)}{g(z;L_\tau)}=\lim_{z\rightarrow 0} \frac z z = 1$$
\en{This yields $\varphi(z;L_\tau)=g(z;L_\tau)$, which proves the Fourier series expansion~(\ref{\en{EN}glgphi}) of $\varphi$.}%
\de{Somit haben wir $\varphi(z;L_\tau)=g(z;L_\tau)$ bewiesen, also die Gleichung~(\ref{\en{EN}glgphi}).}

\en{Finally, we use the definition of $\varphi$ from Prop.~\ref{\en{EN}satzphi} to write $\sigma$ in terms of $\varphi$:}%
\de{Wir müssen nur noch die Definition von $\varphi$ nach $\sigma$ auflösen und erhalten:}
\begin{align*}\sigma(z;L_\tau)&=\exp\lk\frac {\eta_1}{2}\cdot z^2-i\pi z\rk\cdot \varphi(z;L_\tau)\\
&=\exp\lk\frac {\eta_1}{2}\cdot z^2\rk\cdot q_z^{-\frac 1 2}\cdot \frac {q_z-1} {2\pi i}\cdot\prod_{n=1}^\infty\frac{(1-q_\tau^n q_z)(1-q_\tau^n / q_z)}{(1-q_\tau^n)^2} \\
&= \frac 1 {2\pi i}e^{\eta_1\cdot z^2/2}\cdot(q_z^{1/2}-q_z^{-1/2})\cdot\prod_{n=1}^\infty\frac{(1-q_\tau^n q_z)(1-q_\tau^n / q_z)}{(1-q_\tau^n)^2}
\end{align*}
\en{This is the Fourier series representation of the Weierstraß $\sigma$-function from Prop.~\ref{\en{EN}fouriersigma}.}%
\de{Das ist die Fourierdarstellung der Weierstraß'schen $\sigma$-Funktion aus Satz~\ref{\en{EN}fouriersigma}.}
\end{proof}


\begin{theo}\label{\en{EN}fouriertheorem}
\en{Let $q=e^{2\pi i \tau}$. Then $\im(\tau)>0$ yields $|q|<1$ and the following "normalized Eisenstein series" $E_2$, $E_4$ and $E_6$ converge absolutely:}%
\de{Wenn $\im(\tau)>0$ ist, dann gilt für $q=e^{2\pi i \tau}$, dass $|q|<1$ ist und somit die folgenden \glqq normierten Eisensteinreihen\grqq~ $E_2$, $E_4$ und $E_6$ absolut konvergieren:}
\begin{align*}
    E_2(\tau) &:= 1- 24 \sum_{n=1}^\infty n \frac{q^n}{1-q^n}\\
    E_4(\tau) &:= 1+ 240 \sum_{n=1}^\infty n^3 \frac{q^n}{1-q^n}\\
    E_6(\tau) &:= 1- 504 \sum_{n=1}^\infty n^5 \frac{q^n}{1-q^n}
\end{align*}
\en{Using these, we can equivalently redefine the functions $\eta_1(L_\tau)$, $g_2(L_\tau)$ and $g_3(L_\tau)$ which had previously been defined in Def.~\ref{\en{EN}defetak} and Prop.~\ref{\en{EN}dglP}:}%
\de{Mit ihrer Hilfe können wir die in Def.~\ref{\en{EN}defetak} und Satz~\ref{\en{EN}dglP} bereits anders definierten Ausdrücke $\eta_1(L_\tau)$, $g_2(L_\tau)$ und $g_3(L_\tau)$ äquivalent darstellen:}
\begin{align*}
    \eta_1(L_\tau)&=\zeta(z+1;L_\tau)-\zeta(z;L_\tau) = \frac{\pi^2}{3}\cdot E_2(\tau)\\
    g_2(\tau) &= g_2(L_\tau) =60\cdot G_4(L_\tau) =  \frac 4 3 \pi^4 \cdot E_4(\tau)\\
    g_3(\tau) &= g_3(L_\tau) = 140\cdot G_6(L_\tau) = \frac 8 {27} \pi^6 \cdot E_6(\tau)
\end{align*}
\en{And with these new representations of $g_2(L_\tau)$ and $g_3(L_\tau)$ we can also equivalently redefine the discriminant of the lattice $L_\tau$ and its absolute invariant $J$ from Def.~\ref{\en{EN}defijdelta}:}%
\de{Und mit Hilfe der neuen Darstellungen für $g_2(L_\tau)$ und $g_3(L_\tau)$ können wir dann auch die Diskriminante des Gitters $L_\tau$ und die absolute Invariante $J$ aus Def.~\ref{\en{EN}defijdelta} äquivalent darstellen:}
\begin{align*}
    \Delta(\tau) &= \Delta(L_\tau) = \frac{(2\pi)^{12}}{1728} \cdot (E_4(\tau)^3-E_6(\tau)^2)\\
    J(\tau) &= J(L_\tau) = \frac{E_4(\tau)^3}{E_4(\tau)^3-E_6(\tau)^2}
\end{align*}
\en{This expression will not only be called "Klein's absolute invariant of the lattice $L_\tau$", but also "$J$-function", since it associates each $\tau$ from the upper half plane to a complex number.}%
\de{Diesen Ausdruck nennen wir ab jetzt nicht nur \glqq absolute Invariante des Gitters $L_\tau$\grqq, sondern auch \glqq$J$-Funktion\grqq, da sie jedem $\tau$ aus der oberen Halbebene eine komplexe Zahl zuordnet.}
\end{theo}


\begin{proof}
\en{First we calculate the logarithmic derivative of the Fourier series expansion from Prop.~\ref{\en{EN}fouriersigma}.
Then the product yields a sum and we get:}%
\de{Zunächst berechnen wir die logarithmische Ableitung der Fourierentwicklung aus Satz~\ref{\en{EN}fouriersigma}.
Dann geht das Produkt in eine Summe über und wir erhalten:}
\begin{align}
\frac{\sigma'(z;L_\tau)}{\sigma(z;L_\tau)} &= \eta_1\cdot z + \pi i\cdot \frac{e^{\pi i z}+e^{-\pi i z}}{e^{\pi i z}-e^{-\pi i z}}
        + 2\pi i\cdot\sum_{n=1}^\infty \left(\frac{q_\tau^n/q_z}{1-q_\tau^n/q_z}-\frac{q_\tau^n\cdot q_z}{1-q_\tau^n\cdot q_z}\right)\nonumber\\
        &=\eta_1\cdot z + \pi \cdot\frac{\cos(\pi z)}{\sin(\pi z)} + 2\pi i\cdot\sum_{n=1}^\infty \left(\frac{q^n/w}{1-q^n/w}-\frac{q^n\cdot w}{1-q^n\cdot w}\right)\label{\en{EN}koeffi}
\end{align}
\en{where we used $q=q_\tau=e^{2\pi i\tau}$ and $w=q_z=e^{2\pi i z}$ in the last line.
Next we simplify the remaining sum, using the summation formula for geometric series several times:}%
\de{wobei wir $q=q_\tau=e^{2\pi i\tau}$ und $w=q_z=e^{2\pi i z}$ abgekürzt haben.
Als Nächstes vereinfachen wir die verbleibende Summe mit Hilfe der Formel für die geometrische Reihe:}
\begin{align*}
&\phantom{=}~~2\pi i\sum_{n=1}^\infty \left(\frac{q^n/w}{1-q^n/w}-\frac{q^n\cdot w}{1-q^n\cdot w}\right)\\
&= 2\pi i\sum_{n=1}^\infty \sum_{m=1}^\infty\left(\left(q^n/w\right)^m-\left(q^n\cdot w\right)^m\right)= 2\pi i\sum_{n=1}^\infty \sum_{m=1}^\infty(q^m)^n\cdot\left(w^{-m}-w^m\right)\\
&= 2\pi i\sum_{m=1}^\infty \sum_{n=1}^\infty(q^m)^n\cdot\left(w^{-m}-w^m\right) = 2\pi i\sum_{m=1}^\infty \frac{q^m}{1-q^m}\cdot\left(w^{-m}-w^m\right)\\
&= 2\pi i\sum_{m=1}^\infty \frac{q_\tau^m}{1-q_\tau^m}\cdot\left(e^{-2\pi i m z}-e^{2\pi i m z}\right) = 4\pi \sum_{m=1}^\infty \frac{q_\tau^m}{1-q_\tau^m}\cdot\sin(2\pi m z)
\end{align*}
\en{If we put this into equation~(\ref{\en{EN}koeffi}), we get a representation of the $\wp$-function:}%
\de{Wenn wir das jetzt einsetzen, erhalten wir eine Darstellung der $\wp$-Funktion:}
\begin{align*}
\frac{\sigma'(z;L_\tau)}{\sigma(z;L_\tau)} &= \eta_1\cdot z + \pi \cdot\frac{\cos(\pi z)}{\sin(\pi z)} + 4\pi \cdot\sum_{m=1}^\infty \frac{q_\tau^m}{1-q_\tau^m}\cdot\sin(2\pi m z)\quad\left|-\frac{d}{dz}\right.\\
\Longrightarrow\quad \wp(z;L_\tau) &= - \eta_1 + \left(\frac{\pi}{\sin(\pi z)}\right)^2 - 8\pi^2\cdot\sum_{m=1}^\infty \frac{m\cdot q_\tau^m}{1-q_\tau^m}\cdot\cos(2\pi m z)
\end{align*}
\en{Here we use the known Taylor series of $\sin(x)$ and $\cos(x)$ around $x=0$:}%
\de{Nun verwenden wir die Taylorreihen von $\sin(x)$ und $\cos(x)$ bei $x=0$:}
\begin{align*}
    \cos(x)&=1-\frac {x^2}{2!} + \frac {x^4}{4!} + O(x^6)\\
    \Longrightarrow~~ \cos(2\pi m z) &= 1 - 2\pi^2m^2 z^2 + \frac 2 3 \pi^4 m^4 z^4 + O(z^6)\\
    \text{\en{and}\de{und}}\qquad\sin(x)&=x-\frac {x^3}{3!} + \frac {x^5}{5!} - \frac{x^7}{7!}+O(x^9)\\
    \Longrightarrow~~ \frac{\sin(\pi z)}{\pi z} &= 1 - \frac{\pi^2}{6} z^2 + \frac{\pi^4}{120}z^4 - \frac{\pi^6}{5040}z^6 + O(z^8)\\
    \Longrightarrow~~ \left(\frac{\sin(\pi z)}{\pi z}\right)^2 &= 1 - \frac{2\pi^2}{6} z^2 +\left(\frac{2\pi^4}{120}+\frac{\pi^4}{36}\right)z^4 - \left(\frac{2\pi^6}{5040}+\frac{2\pi^6}{6\cdot 120}\right)z^6 + O(z^{8})\\
                    &=1 - \frac{\pi^2}{3} z^2 +\frac{2\pi^4}{45} z^4 - \frac{\pi^6}{315}z^6 + O(z^8)\\
    \Longrightarrow~~ \left(\frac{\pi z}{\sin(\pi z)}\right)^2 &= \left.\left(1-\left(\frac{\pi^2}{3} z^2 - \frac{2\pi^4}{45} z^4 + \frac{\pi^6}{315}z^6 + O(z^8)\right)\right)^{-1}\quad\right|\text{\en{geom.~series}\de{geom.~Reihe}}\\
    &= 1+\left(\frac{\pi^2}{3} z^2 - \frac{2\pi^4}{45} z^4+ \frac{\pi^6}{315}z^6\right)\\
   &~~~~+\left(\frac{\pi^2}{3} z^2 - \frac{2\pi^4}{45} z^4\right)^2+\left(\frac{\pi^2}{3} z^2\right)^3 + O(z^8)\\
    &=\left.1 + \frac{\pi^2}{3} z^2 +\frac{\pi^4}{15} z^4 + \frac{2\pi^6}{189} z^6 + O(z^8)\qquad\right|: z^2\end{align*}\begin{align*}
    \Longrightarrow~~ \left(\frac{\pi}{\sin(\pi z)}\right)^2 &= \frac 1{z^2} + \frac{\pi^2}{3} +\frac{\pi^4}{15} z^2 + \frac{2\pi^6}{189} z^4 + O(z^6)
\end{align*}
\en{This yields the beginning of the Laurent series of $\wp(z;L_\tau)$ around $z=0$:}%
\de{Hiermit erhalten wir den Anfang der Laurentreihe von $\wp(z;L_\tau)$ um $z=0$:}
\begin{align*}
\wp(z;L_\tau) &= - \eta_1 + \left(\frac{\pi}{\sin(\pi z)}\right)^2 - 8\pi^2\cdot\sum_{m=1}^\infty \frac{m\cdot q_\tau^m}{1-q_\tau^m}\cdot\cos(2\pi m z)\\
&= -\eta_1 + \frac{1}{z^2} + \frac{\pi^2}{3} +\frac{\pi^4}{15} z^2 + \frac{2\pi^6}{189} z^4\\
&~~~- 8\pi^2\cdot\sum_{m=1}^\infty \frac{m\cdot q_\tau^m}{1-q_\tau^m}\cdot\left(1 - 2\pi^2m^2 z^2 + \frac 2 3 \pi^4 m^4 z^4\right) + O(z^6)
\end{align*}
\en{In Prop.~\ref{\en{EN}laurentwp} we already calculated the Laurent series of $\wp$:}%
\de{Zur Erinnerung und zum Vergleich hier nochmal die Laurentreihe aus Satz~\ref{\en{EN}laurentwp}:}
$$\wp(z;L)=\frac 1 {z^2} + 3 G_4(L) z^2 + 5 G_6(L) z^4 + \sum_{n=3}^\infty (2n+1)G_{2n+2}(L)\cdot z^{2n}$$
\en{Equating the coefficients of $z^0$, $z^2$ and $z^4$ in these two Laurent series of $\wp$ yields:}%
\de{Ein Vergleich der Koeffizienten vor $z^0$, $z^2$ und $z^4$ in diesen beiden Darstellungen liefert}
\begin{align*}
    0 &= - \eta_1(L_\tau) + \frac {\pi^2}{3} - 8\pi^2\cdot\sum_{m=1}^\infty \frac{m\cdot q_\tau^m}{1-q_\tau^m}\\
    3 G_4(L_\tau) &= \frac{\pi^4}{15} +16\pi^4\cdot\sum_{m=1}^\infty \frac{m^3\cdot q_\tau^m}{1-q_\tau^m}\\
    5 G_6(L_\tau) &= \frac{2\pi^6}{189} - \frac{16}3\pi^6\cdot\sum_{m=1}^\infty \frac{m^5\cdot q_\tau^m}{1-q_\tau^m}
\end{align*}
\en{and we get:}\de{und wir erhalten:}
\begin{align*}
    \eta_1(L_\tau) &= \frac{\pi^2}{3}\left(1-24\sum_{m=1}^\infty \frac{m\cdot q_\tau^m}{1-q_\tau^m}\right)\\
    g_2(L_\tau)&=60 G_4(L_\tau)=\frac{4}{3}\pi^4\left(1+240\sum_{m=1}^\infty \frac{m^3\cdot q_\tau^m}{1-q_\tau^m}\right)\\
    g_3(L_\tau)&=140 G_6(L_\tau)=\frac{8}{27}\pi^6\left(1-504\sum_{m=1}^\infty \frac{m^5\cdot q_\tau^m}{1-q_\tau^m}\right)
\end{align*}

\en{Now we denote the brackets by $E_2(\tau)$, $E_4(\tau)$ and $E_6(\tau)$ and obtain the representations of $\eta_1(\tau)$, $g_2(\tau)$ and $g_3(\tau)$ from Thm.~\ref{\en{EN}fouriertheorem}.}%
\de{Jetzt benennen wir den Inhalt der Klammern mit $E_2(\tau)$, $E_4(\tau)$ und $E_6(\tau)$ und erhalten die in Thm.~\ref{\en{EN}fouriertheorem} genannten Darstellungen für $\eta_1(\tau)$, $g_2(\tau)$ und $g_3(\tau)$.}

\en{Finally, we use these new representations of $g_2$ and $g_3$ in the old Def.~\ref{\en{EN}defijdelta} of $\Delta$ and $J$:}%
\de{Schließlich setzen wir wie angekündigt die neuen Darstellungen von $g_2$ und $g_3$ in die alte Def.~\ref{\en{EN}defijdelta} von $\Delta$ und $J$ ein:}
\begin{align*}
    \Delta(\tau)&=g_2^3(L_\tau)-27g_3^2(L_\tau) = \left(\frac 4 3 \pi^4 \cdot E_4(\tau)\right)^3 - 27 \cdot \left(\frac 8 {27} \pi^6 \cdot E_6(\tau)\right)^2\\
    &=\frac{(2\pi)^{12}}{1728}\cdot\left(E_4(\tau)^3-E_6(\tau)^2\right) \\
    J(\tau)&=\frac{g_2^3(\tau)}{\Delta(\tau)} = \frac{\left(\frac 4 3 \pi^4 \cdot E_4(\tau)\right)^3}{\frac{(2\pi)^{12}}{1728}\cdot\left(E_4(\tau)^3-E_6(\tau)^2\right)} = \frac{E_4(\tau)^3}{E_4(\tau)^3-E_6(\tau)^2}
\end{align*}
\en{Thus we have proven all statements from Theorem~\ref{\en{EN}fouriertheorem}.}%
\de{Somit haben wir alle Aussagen aus Theorem~\ref{\en{EN}fouriertheorem} bewiesen.}
\end{proof}

